\chapter{串联与并联电路}\label{ch:g5:series-parallel-circuits}

一个\omterm{def:g3:first-electric-circuit:bulb}{灯泡}坏了，整串彩灯全都熄灭了——这就是单回路的“魔咒”。
可是在你家厨房里，\omterm{def:g3:first-electric-circuit:bulb}{灯泡}烧坏了，屋里其他的灯依然亮着。
有人把你家的电路接得更巧妙。
今天我们就来学习这种接法：\omterm{def:g5:series-parallel-circuits:parallel}{支路}接法。

\section{两种接线方式}

你已经认识第一种接法：\emph{\omterm{def:g4:series-circuits:series}{串联}}，
所有用电器依次串在同一个回路上，像一串项链上的珠子。下面是第二种接法。

\begin{definition}[并联]\label{def:g5:series-parallel-circuits:parallel}
用电器\emph{并联}\index{并联电路}时，电路分成若干条独立的
\emph{支路}\index{支路}，每条支路只接一个用电器，
各支路先分开、后汇合。道路分叉或汇合的地方
称为
\emph{节点}\index{节点}。每条支路都是从\omterm{def:g3:first-electric-circuit:battery}{电池}的一个\omterm{def:g3:first-electric-circuit:battery}{接线柱}
到另一个\omterm{def:g3:first-electric-circuit:battery}{接线柱}的完整通路。
\end{definition}

\begin{omfigure}
\begin{tikzpicture}[scale=0.8]
  \draw (0,0) to[battery1, l=电池] (0,3.2)
    to[short] (2.2,3.2);
  \draw (2.2,3.2) to[short] (2.2,2.4)
    to[lamp, l=灯泡 1] (5.4,2.4) to[short] (5.4,3.2);
  \draw (2.2,3.2) to[short] (2.2,4.0)
    to[lamp, l=灯泡 2] (5.4,4.0) to[short] (5.4,3.2);
  \draw (5.4,3.2) to[short] (7.0,3.2) to[short] (7.0,0)
    to[short] (0,0);
  \fill (2.2,3.2) circle (2pt);
  \fill (5.4,3.2) circle (2pt);
  \node[above, font=\small] at (3.6,4.95) {两条支路};
\end{tikzpicture}

{\small 两个\omterm{def:g3:first-electric-circuit:bulb}{灯泡}\omterm{def:g5:series-parallel-circuits:parallel}{并联}：电路在一个\omterm{def:g5:series-parallel-circuits:parallel}{节点}处分成两路，
每个\omterm{def:g3:first-electric-circuit:bulb}{灯泡}各占一条\omterm{def:g5:series-parallel-circuits:parallel}{支路}，两条\omterm{def:g5:series-parallel-circuits:parallel}{支路}汇合后，
再回到\omterm{def:g3:first-electric-circuit:battery}{电池}。}
\end{omfigure}

\begin{method}[连接方法]\label{met:g5:series-parallel-circuits:build}
用一节\omterm{def:g3:first-electric-circuit:battery}{电池}、两个\omterm{def:g3:first-electric-circuit:bulb}{灯泡}和导线连接：
\begin{enumerate}
  \item 先把第一个\omterm{def:g3:first-electric-circuit:bulb}{灯泡}单独接在\omterm{def:g3:first-electric-circuit:battery}{电池}上，连成简单回路——
        第一条\omterm{def:g5:series-parallel-circuits:parallel}{支路}；
  \item 再把第二个\omterm{def:g3:first-electric-circuit:bulb}{灯泡}\emph{直接跨接在}第一个\omterm{def:g3:first-electric-circuit:bulb}{灯泡}两端：
        从\omterm{def:g3:first-electric-circuit:bulb}{灯泡} 1 的每一侧各引一根导线到\omterm{def:g3:first-electric-circuit:bulb}{灯泡} 2 的对应一侧——
        第二条\omterm{def:g5:series-parallel-circuits:parallel}{支路}；
  \item 用手指沿每条\omterm{def:g5:series-parallel-circuits:parallel}{支路}走一遍：从 $+$ 出发，
        每条路都应恰好经过一个\omterm{def:g3:first-electric-circuit:bulb}{灯泡}，然后到达 $-$。
\end{enumerate}
两个\omterm{def:g3:first-electric-circuit:bulb}{灯泡}都会亮，
而且每个\omterm{def:g3:first-electric-circuit:bulb}{灯泡}都和单独接一个\omterm{def:g3:first-electric-circuit:bulb}{灯泡}时一样亮。
\end{method}

\section{支路的规律}

\begin{proposition}[并联电路的规律]\label{prop:g5:series-parallel-circuits:rules}
每个用电器各占一条\omterm{def:g5:series-parallel-circuits:parallel}{支路}时：
\begin{enumerate}
  \item \emph{独立性}：一条\omterm{def:g5:series-parallel-circuits:parallel}{支路}断开，只影响
        这一条\omterm{def:g5:series-parallel-circuits:parallel}{支路}——拧下一个\omterm{def:g3:first-electric-circuit:bulb}{灯泡}，另一个\omterm{def:g3:first-electric-circuit:bulb}{灯泡}仍然亮着；
  \item \emph{亮度不变}：每个\omterm{def:g3:first-electric-circuit:bulb}{灯泡}都如同单独接在电路中一样亮，
        \omterm{def:g3:first-electric-circuit:battery}{电池}的“推力”不会在
        各条\omterm{def:g5:series-parallel-circuits:parallel}{支路}之间分享；
  \item \emph{控制位置}：\omterm{def:g5:series-parallel-circuits:parallel}{支路}中的\omterm{ex:g3:first-electric-circuit:switch}{开关}只控制本\omterm{def:g5:series-parallel-circuits:parallel}{支路}；
        装在干路上、分叉之前的\omterm{ex:g3:first-electric-circuit:switch}{开关}控制所有用电器，
        （代价是：\omterm{def:g3:first-electric-circuit:battery}{电池}同时给每条\omterm{def:g5:series-parallel-circuits:parallel}{支路}供电，
        因此耗电更快。）
\end{enumerate}
\end{proposition}

\begin{example}[拧下灯泡的实验再做一次]\label{ex:g5:series-parallel-circuits:test}
对\omterm{def:g5:series-parallel-circuits:parallel}{并联}的那对\omterm{def:g3:first-electric-circuit:bulb}{灯泡}做与\omterm{def:g4:series-circuits:series}{串联}时完全相同的事：
拧下\omterm{def:g3:first-electric-circuit:bulb}{灯泡} 2。\omterm{def:g3:first-electric-circuit:bulb}{灯泡} 1 甚至不会闪一下——它自己的\omterm{def:g5:series-parallel-circuits:parallel}{支路}、
自己完整的通路都未受影响。把\omterm{def:g3:first-electric-circuit:bulb}{灯泡} 2 拧回去，改拧下\omterm{def:g3:first-electric-circuit:bulb}{灯泡} 1，
\omterm{def:g3:first-electric-circuit:bulb}{灯泡} 2 仍然平稳地亮着。单回路的“魔咒”
在\omterm{def:g5:series-parallel-circuits:parallel}{并联电路}中被解除了：每个\omterm{def:g3:first-electric-circuit:bulb}{灯泡}都只管自己的\omterm{def:g5:series-parallel-circuits:parallel}{支路}。
\end{example}

\begin{omfigure}
\begin{tikzpicture}[scale=0.8]
  \draw (0,0) to[battery1] (0,3.2)
    to[cspst, l=总开关] (2.6,3.2);
  \draw (2.6,3.2) to[short] (2.6,2.4)
    to[lamp, l=灯泡 1] (5.8,2.4) to[short] (5.8,3.2);
  \draw (2.6,3.2) to[short] (2.6,4.0)
    to[spst, l=支路开关] (4.2,4.0)
    to[lamp] (5.8,4.0) to[short] (5.8,3.2);
  \draw (5.8,3.2) to[short] (7.4,3.2) to[short] (7.4,0)
    to[short] (0,0);
  \fill (2.6,3.2) circle (2pt);
  \fill (5.8,3.2) circle (2pt);
\end{tikzpicture}

{\small 控制位置：总\omterm{ex:g3:first-electric-circuit:switch}{开关}在分叉之前，同时控制两条\omterm{def:g5:series-parallel-circuits:parallel}{支路}；
\omterm{def:g5:series-parallel-circuits:parallel}{支路}\omterm{ex:g3:first-electric-circuit:switch}{开关}只控制自己的\omterm{def:g3:first-electric-circuit:bulb}{灯泡}——此时它是断开的，
所以\omterm{def:g3:first-electric-circuit:bulb}{灯泡} 1 单独发光。}
\end{omfigure}

\begin{example}[你家是一个并联电路]\label{ex:g5:series-parallel-circuits:house}
你家里的每盏灯、每个插座和每件电器，都各占一条\omterm{def:g5:series-parallel-circuits:parallel}{支路}，
共同组成一个大的\omterm{def:g5:series-parallel-circuits:parallel}{并联电路}。这就是为什么厨房的\omterm{def:g3:first-electric-circuit:bulb}{灯泡}烧坏了，
客厅的灯仍然亮着；为什么每个房间的\omterm{ex:g3:first-electric-circuit:switch}{开关}只控制自己的灯，
以及为什么电表旁保险丝盒里的总\omterm{ex:g3:first-electric-circuit:switch}{开关}，
可以在检修时一次切断全屋的电路。
现代的彩灯也用上了这种接法：它们的\omterm{def:g3:first-electric-circuit:bulb}{灯泡}这样连接，
一个\omterm{def:g3:first-electric-circuit:bulb}{灯泡}坏了，整串彩灯不会再全部熄灭。
\end{example}

\begin{remark}[串联与并联对比]\label{rem:g5:series-parallel-circuits:compare}
断一处：\omterm{def:g4:series-circuits:series}{串联}——全部熄灭；\omterm{def:g5:series-parallel-circuits:parallel}{并联}——只有一条\omterm{def:g5:series-parallel-circuits:parallel}{支路}熄灭。\omterm{def:g3:first-electric-circuit:bulb}{灯泡}增多时：
\omterm{def:g4:series-circuits:series}{串联}——全部变暗；\omterm{def:g5:series-parallel-circuits:parallel}{并联}——全部保持正常亮度，
但\omterm{def:g3:first-electric-circuit:battery}{电池}耗电更快。一个\omterm{ex:g3:first-electric-circuit:switch}{开关}控制全部：\omterm{def:g4:series-circuits:series}{串联}——放在哪里都行；\omterm{def:g5:series-parallel-circuits:parallel}{并联}——只能放在分叉之前。
两种连接都谈不上“更好”。手电筒用\omterm{def:g4:series-circuits:series}{串联}很合适（一个\omterm{def:g3:first-electric-circuit:bulb}{灯泡}，一个\omterm{ex:g3:first-electric-circuit:switch}{开关}），
家中电路则一定是\omterm{def:g5:series-parallel-circuits:parallel}{并联}。电工看任何电路，第一问总是：
电工看任何电路，第一问总是：
电工看任何电路，第一问总是：\emph{道路在何处
分叉？}
\end{remark}

\section{读懂电路}

\begin{method}[沿着路径走]\label{met:g5:series-parallel-circuits:follow}
要预测电路图中任一\omterm{def:g3:first-electric-circuit:bulb}{灯泡}的状态：
\begin{enumerate}
  \item 把手指放在\omterm{def:g3:first-electric-circuit:battery}{电池}的 $+$ 极上，
        沿每条可能的路走到 $-$ 极，
        列出每条路经过的元件；
  \item 若至少有一条未断开的路径经过某\omterm{def:g3:first-electric-circuit:bulb}{灯泡}，
        该\omterm{def:g3:first-electric-circuit:bulb}{灯泡}就会亮；
  \item 同时位于每一条路径上的元件彼此\omterm{def:g4:series-circuits:series}{串联}；
        位于不同分叉上的元件彼此\omterm{def:g5:series-parallel-circuits:parallel}{并联}。
\end{enumerate}
预测之前，先把路径走一遍——
只凭图形的样子去猜，正是电路最容易骗人的地方。
\end{method}

\section{练习}

\begin{exercise}[$\star$]\label{exo:g5:series-parallel-circuits:1}
什么是\omterm{def:g5:series-parallel-circuits:parallel}{支路}？什么是\omterm{def:g5:series-parallel-circuits:parallel}{节点}？在两个\omterm{def:g3:first-electric-circuit:bulb}{灯泡}\omterm{def:g5:series-parallel-circuits:parallel}{并联}的电路中，
有几条\omterm{def:g5:series-parallel-circuits:parallel}{支路}上接有\omterm{def:g3:first-electric-circuit:bulb}{灯泡}？
\end{exercise}

\begin{exercise}[$\star$]\label{exo:g5:series-parallel-circuits:2}
写出\omterm{def:g5:series-parallel-circuits:parallel}{并联电路}的三条规律。
\end{exercise}

\begin{exercise}[$\star$]\label{exo:g5:series-parallel-circuits:3}
两个\omterm{def:g3:first-electric-circuit:bulb}{灯泡}\omterm{def:g5:series-parallel-circuits:parallel}{并联}：\omterm{def:g3:first-electric-circuit:bulb}{灯泡} 1 被拧下。\omterm{def:g3:first-electric-circuit:bulb}{灯泡} 2 会怎样？
\omterm{def:g4:series-circuits:series}{串联}时，\omterm{def:g3:first-electric-circuit:bulb}{灯泡} 2 的情况如何，
两种情况为何不同？
\end{exercise}

\begin{exercise}[$\star$]\label{exo:g5:series-parallel-circuits:4}
两个相同的\omterm{def:g3:first-electric-circuit:bulb}{灯泡}：A 组\omterm{def:g4:series-circuits:series}{串联}，B 组\omterm{def:g5:series-parallel-circuits:parallel}{并联}，
\omterm{def:g3:first-electric-circuit:battery}{电池}相同。哪一组中每个\omterm{def:g3:first-electric-circuit:bulb}{灯泡}更亮？
哪一组的\omterm{def:g3:first-electric-circuit:battery}{电池}耗电更快？
\end{exercise}

\begin{exercise}[$\star$]\label{exo:g5:series-parallel-circuits:5}
在有两个\omterm{ex:g3:first-electric-circuit:switch}{开关}的图中，下列情况下各\omterm{def:g3:first-electric-circuit:bulb}{灯泡}怎样？
两个\omterm{ex:g3:first-electric-circuit:switch}{开关}都闭合？\ 总\omterm{ex:g3:first-electric-circuit:switch}{开关}断开？\ 总\omterm{ex:g3:first-electric-circuit:switch}{开关}闭合，\omterm{def:g5:series-parallel-circuits:parallel}{支路}\omterm{ex:g3:first-electric-circuit:switch}{开关}
断开？
\end{exercise}

\begin{exercise}[$\star$]\label{exo:g5:series-parallel-circuits:6}
家庭电路为什么采用\omterm{def:g5:series-parallel-circuits:parallel}{并联}而非\omterm{def:g4:series-circuits:series}{串联}？
请举出依赖\omterm{def:g5:series-parallel-circuits:parallel}{并联}的两个日常便利。
\end{exercise}

\begin{exercise}[$\star$]\label{exo:g5:series-parallel-circuits:7}
家庭电路中，哪个\omterm{ex:g3:first-electric-circuit:switch}{开关}能同时控制全部用电器？
这个\omterm{ex:g3:first-electric-circuit:switch}{开关}日常叫什么？有什么用途？
\end{exercise}

\begin{exercise}[$\star$]\label{exo:g5:series-parallel-circuits:8}
床头灯及其墙上插座，与天花板灯是\omterm{def:g4:series-circuits:series}{串联}还是
\omterm{def:g5:series-parallel-circuits:parallel}{并联}？你能从日常生活中怎样判断？
\end{exercise}

\begin{exercise}[$\star\star$]\label{exo:g5:series-parallel-circuits:9}
用符号画出走廊的电路图：一节\omterm{def:g3:first-electric-circuit:battery}{电池}、一个\omterm{def:g3:first-electric-circuit:bulb}{灯泡}，
并且还要实现这一便利——\omterm{def:g3:first-electric-circuit:bulb}{灯泡}须能在
走廊\emph{任一}端的\omterm{ex:g3:first-electric-circuit:switch}{开关}\dots\ 请认真尝试；若你的
设计始终失败，请说明困难所在。（真正的双控
需要比我们所用\omterm{ex:g3:first-electric-circuit:switch}{开关}更巧妙的\omterm{ex:g3:first-electric-circuit:switch}{开关}——尝试本身
比答案更有启发。）
\end{exercise}

\begin{exercise}[$\star\star$]\label{exo:g5:series-parallel-circuits:10}
一串现代派对彩灯在一只\omterm{def:g3:first-electric-circuit:bulb}{灯泡}损坏后仍然发光，
——但盒上警告：“拆去的\omterm{def:g3:first-electric-circuit:bulb}{灯泡}不得超过
三只。”请运用\omterm{def:g5:series-parallel-circuits:parallel}{支路}规律，解释这一便利，
并（难度更高地）说明此警告的合理原因。（想想
每条剩下的\omterm{def:g5:series-parallel-circuits:parallel}{支路}从\omterm{def:g3:first-electric-circuit:battery}{电池}得到什么，以及
\omterm{def:g3:first-electric-circuit:battery}{电池}自身的负担。）
\end{exercise}

\begin{exercise}[$\star\star$]\label{exo:g5:series-parallel-circuits:11}
在两个\omterm{ex:g3:first-electric-circuit:switch}{开关}的图中，电工想再加一个\omterm{def:g3:first-electric-circuit:bulb}{灯泡}，
使其\emph{仅}在\omterm{def:g5:series-parallel-circuits:parallel}{支路}\omterm{ex:g3:first-electric-circuit:switch}{开关}闭合时发光，
其余情况下都不亮——即使总\omterm{ex:g3:first-electric-circuit:switch}{开关}闭合。新\omterm{def:g3:first-electric-circuit:bulb}{灯泡}应接在哪一段
路径上？
接好后，它与\omterm{def:g3:first-electric-circuit:bulb}{灯泡} 2 是\omterm{def:g4:series-circuits:series}{串联}还是\omterm{def:g5:series-parallel-circuits:parallel}{并联}？
\end{exercise}
